%% ============================================================
%% Introduction to Cognitive Psychology — Week 8 Study Notes
%% ============================================================
\documentclass[11pt,a4paper]{article}

\usepackage[a4paper, left=1.8cm, right=1.8cm, top=1.3cm, bottom=1.8cm]{geometry}
\usepackage{booktabs}
\usepackage{tabularx}
\usepackage{array}
\usepackage{enumitem}
\usepackage{titlesec}
\usepackage{fancyhdr}
\usepackage{fontenc}
\usepackage{inputenc}
\usepackage{microtype}
\usepackage{parskip}
\usepackage{hyperref}
\usepackage{amsmath}
\usepackage{amssymb}

\hypersetup{hidelinks, colorlinks=false}

%% ── Table helpers ─────────────────────────────────────────
\newcolumntype{L}[1]{>{\raggedright\arraybackslash}p{#1}}

%% ── Section headings ──────────────────────────────────────
\titleformat{\section}{\large\bfseries}{}{0em}{}[\titlerule]
\titleformat{\subsection}{\normalsize\bfseries}{}{0em}{}
\titleformat{\subsubsection}{\normalsize\bfseries}{}{0em}{}

\titlespacing{\section}{0pt}{12pt}{4pt}
\titlespacing{\subsection}{0pt}{8pt}{3pt}
\titlespacing{\subsubsection}{0pt}{6pt}{2pt}

%% ── Page style ────────────────────────────────────────────
\pagestyle{fancy}
\fancyhf{}
\renewcommand{\headrulewidth}{0pt}
\fancyfoot[C]{\small Cognitive Psychology --- Week 8 \textbar\ Page \thepage}

%% ── Misc helpers ──────────────────────────────────────────
\newcommand{\bb}[1]{\textbf{#1}}
\setlength{\parskip}{4pt}

%% ═══════════════════════════════════════════════════════════
\begin{document}
\setlength{\arrayrulewidth}{0.4pt}

\begin{center}
\bfseries\LARGE Introduction to Cognitive Psychology\par
\vspace{0.2cm}
\large Assignment 8 Detailed Solutions
\end{center}
\vspace{0.5cm}

%% ════════════════════════════════════════════════════════════
\section*{ASSIGNMENT 8 --- Full Solutions with Explanations}

\subsubsection*{Question 1}
\textit{Participants doing mental imagery tasks tend to show brain activity in the:}
\medskip

\begin{tabular}{L{0.3cm} L{14.9cm}}
& A.\ Temporal lobe \\
& B.\ Parietal lobe \\
& C.\ Frontal lobe \\
& \textbf{D.\ Occipital lobe \checkmark} \\
\end{tabular}

\vspace{0.3cm}
\noindent\textbf{Ans:} D — Occipital lobe\\[0.2cm]
\textbf{Why this answer:} \\
The \textbf{occipital lobe} is the primary brain region responsible for \textbf{visual processing}. Research in cognitive neuroscience has demonstrated that when participants engage in mental imagery tasks — such as imagining a scene or visualizing an object — they activate many of the same brain areas in the occipital lobe that are used during actual visual perception. This overlap between visual perception and mental imagery supports the idea that imagery is functionally and neurologically linked to perception.
\vspace{0.4cm}

\subsubsection*{Question 2}
\textit{How do people represent and navigate in and through space? This is a question of:}
\medskip

\begin{tabular}{L{0.3cm} L{14.9cm}}
& A.\ Navigational cognition \\
& B.\ Representational cognition \\
& \textbf{C.\ Spatial cognition \checkmark} \\
& D.\ Visual cognition \\
\end{tabular}

\vspace{0.3cm}
\noindent\textbf{Ans:} C — Spatial cognition\\[0.2cm]
\textbf{Why this answer:} \\
\textbf{Spatial cognition} is the branch of cognitive psychology concerned with how organisms \textbf{acquire, organize, and use knowledge about their spatial environment}. It encompasses how people form mental representations of space (cognitive maps), navigate through environments, estimate distances and directions, and reason about spatial relationships. While ``visual cognition'' deals broadly with visual processing, and ``navigational cognition'' is not a standard term, \textit{spatial cognition} is the precise term used in the field.
\vspace{0.4cm}

\subsubsection*{Question 3}
\textit{Images can prime the visual pathway, making it easier to detect a faint stimulus. This is an example of:}
\medskip

\begin{tabular}{L{0.3cm} L{14.9cm}}
& \textbf{A.\ Perceptual equivalence \checkmark} \\
& B.\ Spatial equivalence \\
& C.\ Structural equivalence \\
& D.\ Transformational equivalence \\
\end{tabular}

\vspace{0.3cm}
\noindent\textbf{Ans:} A — Perceptual equivalence\\[0.2cm]
\textbf{Why this answer:} \\
Finke's principle of \textbf{perceptual equivalence} states that mental images activate the \textbf{same perceptual mechanisms} as actual visual stimuli. When a mental image ``primes'' the visual pathway and facilitates the detection of a faint stimulus, it demonstrates that imagery and perception share common neural and functional pathways. This principle is foundational evidence for the claims that imagery is not purely abstract but engages genuine perceptual processing.
\vspace{0.4cm}

\subsubsection*{Question 4}
\textit{Which of the following questions would take you the LONGEST TIME to answer?}
\medskip

\begin{tabular}{L{0.3cm} L{14.9cm}}
& A.\ Which is faster, a turtle or a sports car? \\
& B.\ Which is faster, a turtle or a cheetah? \\
& \textbf{C.\ Which is faster, a turtle or a caterpillar? \checkmark} \\
& D.\ Which is faster, a caterpillar or a cat? \\
\end{tabular}

\vspace{0.3cm}
\noindent\textbf{Ans:} C — Which is faster, a turtle or a caterpillar?\\[0.2cm]
\textbf{Why this answer:} \\
This question relates to the \textbf{symbolic distance effect} in mental imagery. When two items being compared are \textbf{very similar} on the dimension of comparison (in this case, speed), it takes \textbf{longer} to make a judgment. A turtle and a caterpillar are both very slow creatures, making the speed difference between them small and the comparison more difficult. In contrast, comparisons like turtle vs.\ sports car or turtle vs.\ cheetah involve large differences that are resolved quickly.
\vspace{0.4cm}

\subsubsection*{Question 5}
\textit{The relational-organizational hypothesis is supported by:}
\medskip

\begin{tabular}{L{0.3cm} L{14.9cm}}
& A.\ The effectiveness of the method of loci \\
& B.\ The fact that concrete words are recalled better than abstract words \\
& \textbf{C.\ The fact that non-interactive images do not facilitate recall whereas interactive images do facilitate recall \checkmark} \\
& D.\ The effectiveness of the pegword method \\
\end{tabular}

\vspace{0.3cm}
\noindent\textbf{Ans:} C — The fact that non-interactive images do not facilitate recall whereas interactive images do facilitate recall\\[0.2cm]
\textbf{Why this answer:} \\
The \textbf{relational-organizational hypothesis} proposes that imagery aids memory primarily because it helps organize items into \textbf{meaningful relational structures}. The critical evidence is that simply forming separate, non-interactive images of items does \textbf{not} improve recall — what matters is creating \textbf{interactive images} where items are linked in a meaningful relationship. This demonstrates that it is the relational organization, not mere visual representation, that enhances memory performance.
\vspace{0.4cm}


\subsubsection*{Question 6}
\textit{Paivio's \_\_\_\_ hypothesis argues that long-term memory contains two separate systems that represent information in verbal and visual forms, respectively.}
\medskip

\begin{tabular}{L{0.3cm} L{14.9cm}}
& A.\ Picture-word \\
& \textbf{B.\ Dual code \checkmark} \\
& C.\ Visuo-verbal \\
& D.\ Relational-Organizational \\
\end{tabular}

\vspace{0.3cm}
\noindent\textbf{Ans:} B — Dual code\\[0.2cm]
\textbf{Why this answer:} \\
Allan Paivio's \textbf{dual-code hypothesis} (also known as dual-coding theory) proposes that long-term memory consists of \textbf{two separate but interconnected coding systems}: a \textbf{verbal system} (logogens) for processing and storing linguistic information, and a \textbf{visual/imaginal system} (imagens) for processing and storing visual-spatial information. Concrete words benefit from dual coding (both systems), explaining their recall advantage over abstract words that are primarily encoded verbally.
\vspace{0.4cm}
\newpage
\subsubsection*{Question 7}
\textit{Which of the following methods would most help you to remember the word pair ``elephant-cigar''?}
\medskip

\begin{tabular}{L{0.3cm} L{14.9cm}}
& A.\ Repeating the word elephant over and over while visualizing a cigar \\
& \textbf{B.\ Visualizing an elephant smoking a cigar \checkmark} \\
& C.\ Repeating both words over and over to yourself \\
& D.\ Visualizing an elephant and a cigar, not touching each other \\
\end{tabular}

\vspace{0.3cm}
\noindent\textbf{Ans:} B — Visualizing an elephant smoking a cigar\\[0.2cm]
\textbf{Why this answer:} \\
Research on \textbf{interactive imagery} consistently shows that forming a \textbf{vivid, interactive mental image} linking two items together produces \textbf{superior recall} compared to all other strategies. Visualizing an elephant \textit{smoking} a cigar creates a meaningful interaction between the items. Simply visualizing them separately (option D), or using rote repetition (options A and C), does not create the kind of relational encoding that makes the association memorable. Interactive imagery leverages the relational-organizational benefits of mental imagery.
\vspace{0.4cm}

\subsubsection*{Question 8}
\textit{Your mental map of your campus is probably:}
\medskip

\begin{tabular}{L{0.3cm} L{14.9cm}}
& A.\ Larger than reality \\
& B.\ Smaller than reality \\
& \textbf{C.\ More regular, with more straight lines and right angles, than reality \checkmark} \\
& D.\ Less regular, with fewer straight lines and right angles, than reality \\
\end{tabular}

\vspace{0.3cm}
\noindent\textbf{Ans:} C — More regular, with more straight lines and right angles, than reality\\[0.2cm]
\textbf{Why this answer:} \\
Research on \textbf{cognitive maps} has consistently shown that people's mental representations of spatial environments are \textbf{systematically distorted}. One of the most common distortions is \textbf{regularization} — people tend to straighten curved roads, make intersections into right angles, and align landmarks more neatly than they actually are. This reflects a tendency to impose more geometric regularity and simplicity on spatial representations, making them easier to store and process but less accurate in detail.
\vspace{0.4cm}

\subsubsection*{Question 9}
\textit{Finke's principle of \_\_\_\_\_ states that mental imagery allows us to retrieve information that was not intentionally stored.}
\medskip

\begin{tabular}{L{0.3cm} L{14.9cm}}
& A.\ Perceptual equivalence \\
& B.\ Transformational equivalence \\
& \textbf{C.\ Implicit encoding \checkmark} \\
& D.\ Structural encoding \\
\end{tabular}

\vspace{0.3cm}
\noindent\textbf{Ans:} C — Implicit encoding\\[0.2cm]
\textbf{Why this answer:} \\
Finke's principle of \textbf{implicit encoding} holds that mental imagery enables people to \textbf{discover or retrieve information that was never deliberately encoded}. For example, after forming a mental image of a familiar object, a person may notice details or spatial relationships they had never consciously attended to before. This principle demonstrates the power of imagery as a cognitive tool — it allows access to information that is embedded in perceptual experience but not explicitly stored or labeled in memory.
\vspace{0.4cm}
\newpage
\subsubsection*{Question 10}
\textit{A propositional representation is thought to be \_\_\_\_\_\_ in nature.}
\medskip

\begin{tabular}{L{0.3cm} L{14.9cm}}
& A.\ Verbal \\
& B.\ Visual \\
& C.\ Both verbal and visual \\
& \textbf{D.\ Neither verbal nor visual \checkmark} \\
\end{tabular}

\vspace{0.3cm}
\noindent\textbf{Ans:} D — Neither verbal nor visual\\[0.2cm]
\textbf{Why this answer:} \\
A \textbf{propositional representation} is an \textbf{abstract, amodal} form of mental representation that captures the meaning of information without being tied to any specific sensory modality. Propositions are \textbf{language-like} in their logical structure (consisting of predicates and arguments), but they are \textbf{not} themselves verbal or visual. They are considered the underlying ``mentalese'' — a level of representation that is deeper and more fundamental than either words or images. This is a key distinction in the imagery debate between analog (image-based) and propositional theories of mental representation.
\vspace{0.4cm}

\medskip
\subsection*{Assignment Quick Answer Summary}

\begin{center}
\renewcommand{\arraystretch}{1.4}
\begin{tabular}{|L{0.8cm}|L{5.5cm}|L{8.9cm}|}
\hline
\bb{Q\#} & \bb{Topic} & \bb{Correct Answer} \\
\hline
1 & Mental Imagery Brain Activity & D — Occipital lobe \\
2 & Spatial Navigation Question & C — Spatial cognition \\
3 & Image Priming Visual Pathway & A — Perceptual equivalence \\
4 & Longest Comparison Time & C — Turtle vs.\ caterpillar \\
5 & Relational-Organizational Hypothesis & C — Interactive images facilitate recall \\
6 & Paivio's Memory Systems & B — Dual code \\
7 & Best Method for Word Pair Memory & B — Visualizing elephant smoking a cigar \\
8 & Mental Map Distortion & C — More regular with straight lines/right angles \\
9 & Finke's Principle of Retrieval & C — Implicit encoding \\
10 & Propositional Representation Nature & D — Neither verbal nor visual \\
\hline
\end{tabular}
\end{center}

\medskip
\end{document}
